\chapter{Vectors in Two Dimensions}

\chapterintro{A vector is a displacement, a velocity, a force, or a direction, all stored as the same pair of numbers. Adding them moves things, normalizing them keeps speeds honest, and the dot product answers the questions a game asks most: how far apart, how aligned, and is it in front of me.}

\section{Core Ideas}

\paragraph{Vectors and components} A vector $\mathbf{v} = \vv{x, y}$ has magnitude and direction but no position. The vector from point $A$ to point $B$ is $B - A$, component by component. A point's position vector is the vector from the origin to it.

\paragraph{Addition and scaling} $\vv{a, b} + \vv{c, d} = \vv{a + c, b + d}$, pictured tip to tail. $k\vv{a, b} = \vv{ka, kb}$ scales the length by $\abs{k}$ and reverses direction when $k < 0$. Subtraction $\mathbf{u} - \mathbf{v}$ is the vector from the tip of $\mathbf{v}$ to the tip of $\mathbf{u}$. Position plus velocity times time is this arithmetic.

\paragraph{Magnitude and unit vectors} $\abs{\vv{x, y}} = \sqrt{x^2 + y^2}$ by Chapter~3. The unit vector in the direction of $\mathbf{v}$ is $\mathbf{v}/\abs{\mathbf{v}}$, which has length 1; multiplying it by a speed gives a velocity with that exact speed. The zero vector has no direction and cannot be normalized. The direction angle is $\operatorname{atan2}(y, x)$.

\paragraph{The dot product} Defined by components and equal to a product of lengths and the cosine of the angle between the vectors. Its sign says whether the angle is acute (positive), right (zero) or obtuse (negative), with no trigonometry needed.

\begin{keybox}{Dot product}
\[
\mathbf{a} \cdot \mathbf{b} = a_x b_x + a_y b_y = \abs{\mathbf{a}}\,\abs{\mathbf{b}}\cos\theta, \qquad
\cos\theta = \frac{\mathbf{a} \cdot \mathbf{b}}{\abs{\mathbf{a}}\,\abs{\mathbf{b}}}
\]
\end{keybox}

\paragraph{Projection} The projection of $\mathbf{a}$ onto $\mathbf{b}$ is the part of $\mathbf{a}$ that points along $\mathbf{b}$: $\proj_{\mathbf{b}}\mathbf{a} = \dfrac{\mathbf{a} \cdot \mathbf{b}}{\abs{\mathbf{b}}^2}\,\mathbf{b}$ (Figure~\ref{fig:proj}). Its length, the scalar projection, is $\mathbf{a} \cdot \mathbf{b} / \abs{\mathbf{b}}$. Sliding along a wall keeps only the projection of the velocity onto the wall.

\begin{figure}[H]
\centering
\begin{tikzpicture}[scale=0.9,every node/.style={font=\small}]
\draw[hairline] (0,0) -- (5.6,0);
\draw[navy,thick,->] (0,0) -- (5,0) node[below,text=navy] {$\mathbf{b}$};
\draw[navy,thick,->] (0,0) -- (3,2.2) node[above,text=navy] {$\mathbf{a}$};
\draw[navy!60,dashed] (3,2.2) -- (3,0);
\draw[navy!60,very thick,->] (0,-0.25) -- (3,-0.25) node[midway,below,text=slate] {$\proj_{\mathbf{b}}\mathbf{a}$};
\draw[navy!60] (2.75,0) -- (2.75,0.25) -- (3,0.25);
\end{tikzpicture}
\caption{The projection of $\mathbf{a}$ onto $\mathbf{b}$ is the shadow $\mathbf{a}$ casts on the line of $\mathbf{b}$.}
\label{fig:proj}
\end{figure}

\paragraph{Perpendicular vectors} In two dimensions, $\vv{-y, x}$ is perpendicular to $\vv{x, y}$ and has the same length; it is the counterclockwise $90\degs$ rotation from Chapter~5. The dot product of perpendicular vectors is zero.

\paragraph{Reflection} A velocity $\mathbf{v}$ bouncing off a surface with unit normal $\mathbf{n}$ becomes $\mathbf{v} - 2(\mathbf{v} \cdot \mathbf{n})\mathbf{n}$: the component along the normal is reversed and the rest is kept.

\section{Worked Examples}

\begin{example}
Let $\mathbf{a} = \vv{3, -1}$ and $\mathbf{b} = \vv{-2, 4}$. Find $\mathbf{a} + \mathbf{b}$, $2\mathbf{a} - \mathbf{b}$ and $\abs{\mathbf{a}}$.
\solution $\mathbf{a} + \mathbf{b} = \vv{1, 3}$. $2\mathbf{a} - \mathbf{b} = \vv{6, -2} - \vv{-2, 4} = \vv{8, -6}$. $\abs{\mathbf{a}} = \sqrt{9 + 1} = \sqrt{10}$.
\end{example}

\begin{example}
Find the angle between $\vv{1, 2}$ and $\vv{3, 1}$.
\solution The dot product is $3 + 2 = 5$. The lengths are $\sqrt{5}$ and $\sqrt{10}$, so $\cos\theta = 5/\sqrt{50} = 1/\sqrt{2}$ and $\theta = 45\degs$.
\end{example}

\begin{example}
Project $\vv{3, 4}$ onto $\vv{1, 1}$.
\solution The dot product is $7$ and $\abs{\vv{1, 1}}^2 = 2$, so the projection is $\tfrac{7}{2}\vv{1, 1} = \vv{3.5, 3.5}$.
\end{example}

\section{Exercises}

\begin{multicols}{2}
\begin{exercises}
\item For $\mathbf{u} = \vv{2, -3}$ and $\mathbf{v} = \vv{-1, 5}$, find $\mathbf{u} + \mathbf{v}$, $\mathbf{u} - \mathbf{v}$, $3\mathbf{u}$ and $\abs{\mathbf{u}}$.
\item Find the displacement vector from $P(1, 4)$ to $Q(-2, 8)$, its magnitude, and the unit vector in its direction.
\item Find the unit vector in the direction of $\vv{6, 8}$, and the vector of length $15$ in that direction.
\item Compute $\vv{3, 4} \cdot \vv{4, -3}$, $\vv{1, 2} \cdot \vv{2, 1}$ and $\vv{-1, 2} \cdot \vv{3, 1}$, and state what each sign says about the angle.
\item Find the angle between $\vv{1, 0}$ and $\vv{1, \sqrt{3}}$.
\item Find the angle between $\vv{2, 2}$ and $\vv{-1, 1}$.
\item Project $\vv{5, 3}$ onto $\vv{2, 0}$ and give the scalar projection. Then project $\vv{2, 6}$ onto $\vv{1, 1}$.
\item Give a vector perpendicular to $\vv{3, 7}$ and verify it with the dot product.
\item Find the direction angle of $\vv{-4, -4}$.
\item A character at $(10, 2)$ faces direction $\vv{1, 0}$ and an enemy stands at $(14, 5)$. Find the vector to the enemy, its dot product with the facing direction, and the angle between them. Is the enemy inside a view cone of half-angle $45\degs$?
\item A ball with velocity $\vv{3, -4}$ hits a horizontal floor with unit normal $\vv{0, 1}$. Find its velocity after the bounce.
\item \textbf{(proof)} Using components, prove that $\mathbf{a} \cdot (\mathbf{b} + \mathbf{c}) = \mathbf{a} \cdot \mathbf{b} + \mathbf{a} \cdot \mathbf{c}$, and that $\vv{-y, x}$ is perpendicular to $\vv{x, y}$ with the same length.
\end{exercises}
\end{multicols}

\section{Programming Exercises}

\begin{codeexercises}
\item Write \lstinline!add!, \lstinline!sub!, \lstinline!scale!, \lstinline!magnitude!, \lstinline!normalize! (raising on the zero vector) and \lstinline!dot! on coordinate tuples. Test with Exercises~10.1 to 10.4.
\item Write \lstinline!angle_between(a, b)! in degrees using \lstinline!math.acos!, clamping the cosine to $[-1, 1]$ first. Test with Exercises~10.5 and 10.6, then explain in a comment why the clamp is needed when \lstinline!a! and \lstinline!b! are parallel.
\item Write \lstinline!project(a, b)! and \lstinline!perp(v)!. Test with Exercises~10.7 and 10.8, and assert that \lstinline!dot(perp(v), v)! is zero for 100 random vectors.
\item Write \lstinline!in_view_cone(pos, facing, target, half_angle_deg)!. Test with Exercise~10.10, then move the target around a circle of radius 5 about the character and print which angles are visible.
\item Write \lstinline!reflect(v, n)! and simulate a ball inside a box for 500 steps, bouncing off all four walls. Assert that the speed never changes, and explain which of the chapter's facts guarantees that.
\end{codeexercises}
